\section{Рабочий проект}
\subsection{Классы, используемые при разработке сайта}

Можно выделить следующий список классов и их методов, использованных при разработке web-приложения (таблицы \ref{tab:crypto_api} - \ref{tab:additional_api}).

\subsubsection{Модуль работы с сессиями}

\begin{xltabular}{\textwidth}{|l|l|p{1.7cm}|X|}
	\caption{Функции модуля session.py}\label{tab:crypto_api} \\ \hline
	\centrow Функция & \centrow Тип & \centrow Метод & \centrow Описание \\ \hline
	\thead{1} & \thead{2} & \centrow 3 & \centrow 4 \\ \hline
	\endfirsthead
	\continuecaption{Продолжение таблицы \ref{tab:crypto_api}}
	\thead{1} & \thead{2} & \centrow 3 & \centrow 4 \\ \hline
	\finishhead
	create\_session() & Генерация & - & Создает уникальный идентификатор сессии и 128-битный ключ. Возвращает кортеж (session\_id, key) \\ \hline
	get\_key(session\_id) & Получение & GET & Извлекает ключ шифрования из БД по идентификатору сессии \\ \hline
	delete\_session(session\_id) & Удаление & DELETE & Удаляет сессию и связанный ключ из БД \\ \hline
\end{xltabular}

\subsubsection{Процесс шифрования/дешифрования}

\begin{xltabular}{\textwidth}{|l|X|}
	\caption{Последовательность работы с шифрованием}\label{tab:crypto_flow} \\ \hline
	\centrow Этап & \centrow Описание \\ \hline
	\thead{1} & \thead{2} \\ \hline
	\endfirsthead
	\continuecaption{Продолжение таблицы \ref{tab:crypto_flow}}
	\finishhead
	1. Инициализация & При входе пользователя создается сессия (create\_session()), session\_id и ключ сохраняются в sessionStorage \\ \hline
	2. Шифрование (сервер) & Сообщение дополняется до размера блока (pad()), шифруется AES-128 (ECB режим) и кодируется в base64 \\ \hline
	3. Передача & Зашифрованное сообщение передается клиенту через API /get\_private\_messages \\ \hline
	4. Дешифрование (клиент) & Клиент использует сохраненный ключ для дешифрования через CryptoJS.AES.decrypt() \\ \hline
	5. Завершение & При выходе сессия удаляется (delete\_session()) \\ \hline
\end{xltabular}

\textbf{Особенности реализации:}
\begin{itemize}
	\item используется режим ECB как наиболее простой для реализации;
	\item ключ сессии генерируется случайно (get\_random\_bytes());
	\item сообщения дополняются по стандарту PKCS7;
	\item все ключи хранятся только в БД и sessionStorage.
\end{itemize}

\subsubsection{Аутентификация и пользователи}

\begin{xltabular}{\textwidth}{|l|l|p{1.7cm}|X|}
	\caption{API для работы с пользователями}\label{tab:users_api} \\ \hline
	\centrow Поле & \centrow Тип & \centrow Метод & \centrow Описание \\ \hline
	\thead{1} & \thead{2} & \centrow 3 & \centrow 4 \\ \hline
	\endfirsthead
	\continuecaption{Продолжение таблицы \ref{tab:users_api}}
	\thead{1} & \thead{2} & \centrow 3 & \centrow 4 \\ \hline
	\finishhead
	/register & Регистрация & POST & Создание нового пользователя. Параметры: username, password \\ \hline 
	/login & Аутентификация & POST & Вход в систему. Параметры: username, password \\ \hline 
	/get\_user\_id & Информация & GET & Получение данных текущего пользователя \\ \hline 
\end{xltabular}

\subsubsection{Работа с сообщениями}

\begin{xltabular}{\textwidth}{|l|l|p{1.7cm}|X|}
	\caption{API для работы с сообщениями}\label{tab:messages_api} \\ \hline
	\centrow Поле & \centrow Тип & \centrow Метод & \centrow Описание \\ \hline
	\thead{1} & \thead{2} & \centrow 3 & \centrow 4 \\ \hline
	\endfirsthead
	\continuecaption{Продолжение таблицы \ref{tab:messages_api}}
	\thead{1} & \thead{2} & \centrow 3 & \centrow 4 \\ \hline
	\finishhead
	/get\_messages & Получение & GET & Запрос новых сообщений. Параметры: timestamp \\ \hline 
	/send\_message & Отправка & POST & Отправка сообщения. Параметры: message, files[] \\ \hline 
	/delete\_message/\{id\} & Удаление & DELETE & Удаление сообщения по ID \\ \hline 
	/edit\_message/\{id\} & Редактирование & PUT & Изменение сообщения. Параметры: message \\ \hline 
	/check\_messages & Проверка & GET & Проверка существования сообщений. Параметры: ids[] \\ \hline 
	/check\_edited\_messages & Проверка & GET & Поиск измененных сообщений. Параметры: last\_timestamp \\ \hline 
\end{xltabular}

\subsubsection{Групповые чаты}

\begin{xltabular}{\textwidth}{|l|l|p{1.7cm}|X|}
	\caption{API для работы с группами}\label{tab:groups_api} \\ \hline
	\centrow Поле & \centrow Тип & \centrow Метод & \centrow Описание \\ \hline
	\thead{1} & \thead{2} & \centrow 3 & \centrow 4 \\ \hline
	\endfirsthead
	\continuecaption{Продолжение таблицы \ref{tab:groups_api}}
	\thead{1} & \thead{2} & \centrow 3 & \centrow 4 \\ \hline
	\finishhead
	/create\_group & Создание & POST & Создание новой группы. Параметры: name \\ \hline 
	/add\_to\_group & Добавление & POST & Приглашение пользователя. Параметры: group\_id, username \\ \hline 
	/get\_groups & Список & GET & Получение групп пользователя \\ \hline 
	/get\_group\_messages & Сообщения & GET & Получение сообщений группы. Параметры: group\_id \\ \hline 
	/get\_group\_members & Участники & GET & Получение списка участников. Параметры: group\_id \\ \hline 
	/leave\_group & Выход & POST & Покидание группы. Параметры: group\_id \\ \hline 
	/change\_member\_role & Права & POST & Изменение роли. Параметры: group\_id, username, role \\ \hline 
	/rename\_group & Переименование & POST & Изменение названия. Параметры: group\_id, new\_name \\ \hline 
	/remove\_from\_group & Исключение & POST & Удаление участника. Параметры: group\_id, username \\ \hline 
		 	 \end{xltabular}

\subsubsection{Личные сообщения}

\begin{xltabular}{\textwidth}{|l|l|p{1.7cm}|X|}
	\caption{API для личных сообщений}\label{tab:private_api} \\ \hline
	\centrow Поле & \centrow Тип & \centrow Метод & \centrow Описание \\ \hline
	\thead{1} & \thead{2} & \centrow 3 & \centrow 4 \\ \hline
	\endfirsthead
	\continuecaption{Продолжение таблицы \ref{tab:private_api}}
	\thead{1} & \thead{2} & \centrow 3 & \centrow 4 \\ \hline
	\finishhead
	/send\_private\_message & Отправка & POST & Отправка ЛС. Параметры: receiver, message \\ \hline 
	/get\_private\_messages & Получение & GET & Получение ЛС. Параметры: user, timestamp \\ \hline 
	/get\_private\_chats & Список & GET & Получение чатов пользователя \\ \hline 
	/search\_users & Поиск & GET & Поиск пользователей. Параметры: q \\ \hline 
\end{xltabular}

\subsubsection{Дополнительные API}

\begin{xltabular}{\textwidth}{|l|l|p{1.7cm}|X|}
	\caption{Дополнительные API эндпоинты}\label{tab:additional_api} \\ \hline
	\centrow Поле & \centrow Тип & \centrow Метод & \centrow Описание \\ \hline
	\thead{1} & \thead{2} & \centrow 3 & \centrow 4 \\ \hline
	\endfirsthead
	\continuecaption{Продолжение таблицы \ref{tab:additional_api}}
	\thead{1} & \thead{2} & \centrow 3 & \centrow 4 \\ \hline
	\finishhead
	/get\_group\_name & Информация & GET & Получение названия группы по ID \\ \hline 
	/send\_system\_message & Системное & POST & Отправка системных сообщений \\ \hline 
	/check\_groups\_updates & Проверка & GET & Проверка обновлений в группах \\ \hline 
	/check\_private\_chats\_updates & Проверка & GET & Проверка новых ЛС \\ \hline 
\end{xltabular}

\subsubsection{Описание классов модуля auth.py}

Модуль auth.py отвечает за регистрацию, вход, выход из системы и удаление сессий. Именно он запускает пользователя в систему и следит за тем, чтобы он был авторизован.

\begin{xltabular}{\textwidth}{|l|X|}
	\caption{Классы модуля auth.py} \label{tab:auth_classes} \\ \hline
	Название класса & Описание класса \\ \hline
	\endfirsthead
	\continuecaption{Продолжение таблицы \ref{tab:auth_classes}}
	Название класса & Описание класса \\ \hline
	\finishhead
	RegisterView & Класс, реализующий отображение формы регистрации и обработку POST-запроса для создания нового пользователя. При успешной регистрации возвращается сообщение об успехе, а в случае возникновения ошибки выводится соответствующее уведомление. \\ \hline
	LoginView & Обеспечивает отображение формы входа и обработку аутентификации пользователя. При успешном логине создается сессия, передаются параметры сессии (session\_id, session\_key) и происходит перенаправление на главную страницу. \\ \hline
	LogoutView & Обрабатывает запрос на выход из системы, удаляет сессию, очищает cookies и перенаправляет пользователя на страницу логина. \\ \hline
	DeleteSessionView & Позволяет удалить сессию по переданному идентификатору. Используется для очистки сессии на сервере. \\ \hline
\end{xltabular}

\subsection{Пакет Views}

Пакет views отвечает за реализацию обработчиков HTTP-запросов, связывая маршруты и бизнес-логику с пользовательским интерфейсом. Он содержит классы, каждый из которых соответствует определённому функциональному блоку системы, таким как регистрация пользователей, отправка и получение сообщений, управление группами, работа с личными чатами, поиск и отображение ошибок. Каждый модуль внутри пакета решает отдельную прикладную задачу и реализует определённую категорию интерфейсного поведения — от авторизации до взаимодействия с базой данных сообщений. Пакет views фактически формирует ядро взаимодействия между клиентской частью мессенджера и сервером, обрабатывая запросы, проверяя права доступа, формируя ответы в нужном формате и обеспечивая стабильную и логически разделённую архитектуру веб-приложения. Благодаря чёткому разграничению ответственности между модулями, этот пакет обеспечивает поддерживаемость, читаемость и масштабируемость кода.

\subsubsection{Описание классов модуля base.py}

Модуль base.py содержит базовые абстракции (View, TemplateView), от которых наследуются другие обработчики, а также класс IndexView, рендерящий главную страницу.

\begin{xltabular}{\textwidth}{|l|X|}
	\caption{Классы модуля base.py} \label{tab:base_classes} \\ \hline
	Название класса & Описание класса \\ \hline
	\endfirsthead
	\continuecaption{Продолжение таблицы \ref{tab:base_classes}}
	Название класса & Описание класса \\ \hline
	\finishhead
	View & Базовый класс для всех обработчиков. Реализует метод \textit{response()}, который по умолчанию пытается загрузить файл с диска и возвращает его с установленным MIME-типом. \\ \hline
	TemplateView & Наследник класса View, предназначенный для работы с шаблонами. Читает файл шаблона и возвращает его содержимое с корректным MIME-типом. \\ \hline
	IndexView & Конкретная реализация TemplateView для отображения главной страницы приложения корпоративного мессенджера. \\ \hline
\end{xltabular}

\subsubsection{Описание классов модуля message.py}

Модуль message.py - главный мотор общения: получение, отправка, редактирование, удаление сообщений, работа с вложениями, системными событиями, проверками изменений.

\begin{xltabular}{\textwidth}{|l|X|}
	\caption{Классы модуля message.py} \label{tab:message_classes} \\ \hline
	Название класса & Описание класса \\ \hline
	\endfirsthead
	\continuecaption{Продолжение таблицы \ref{tab:message_classes}}
	Название класса & Описание класса \\ \hline
	\finishhead
	GetMessageView & Обеспечивает получение сообщений общего чата, отбирая те, что новее указанного timestamp. \\ \hline
	SendMessageView & Обрабатывает отправку сообщений. Поддерживает различные типы сообщений (общие, групповые, личные) и обработку вложенных файлов. \\ \hline
	DeleteMessageView & Обрабатывает запрос на удаление сообщения. Проверяет, что пользователь имеет право удалить данное сообщение, и удаляет его вместе с вложениями. \\ \hline
	EditMessageView & Позволяет редактировать текст сообщения. Производится проверка прав пользователя, после чего обновляется запись в базе данных с новым текстом и timestamp. \\ \hline
	GetGroupMessagesView & Получает сообщения группового чата для заданной группы, фильтруя их по timestamp. \\ \hline
	GetPrivateMessagesView & Обеспечивает получение личных сообщений между двумя пользователями с учетом фильтрации по времени. \\ \hline
	SendPrivateMessageView & Обрабатывает отправку личного сообщения между пользователями. \\ \hline
	SendSystemMessageView & Позволяет отправлять системные сообщения (например, уведомления об изменениях в группе) от имени системы. \\ \hline
	CheckMessagesView & Проверяет наличие определённых сообщений (например, для синхронизации с клиентом) и возвращает список существующих идентификаторов сообщений. \\ \hline
	CheckEditedMessagesView & Проверяет, какие сообщения были изменены после указанного timestamp, и возвращает сведения об отредактированных сообщениях. \\ \hline
\end{xltabular}

\subsubsection{Описание классов модуля p\_chat.py}

Модуль p\_chat.py обслуживает личные диалоги один-на-один.

\begin{xltabular}{\textwidth}{|l|X|}
	\caption{Классы модуля p\_chat.py} \label{tab:p_chat_classes} \\ \hline
	Название класса & Описание класса \\ \hline
	\endfirsthead
	\continuecaption{Продолжение таблицы \ref{tab:p_chat_classes}}
	Название класса & Описание класса \\ \hline
	\finishhead
	GetPrivateChatsView & Получает список личных чатов для текущего пользователя, формируя список собеседников и определяя время последней активности в каждом чате. \\ \hline
	CheckPrivateChatsUpdatesView & Проверяет наличие новых сообщений в личных чатах, что позволяет клиенту определить, требуется ли обновление списка чатов. \\ \hline
\end{xltabular}

\subsubsection{Описание классов модуля search.py}

Модуль search.py отвечает за поиск пользователей и сообщений.

\begin{xltabular}{\textwidth}{|l|X|}
	\caption{Классы модуля search.py} \label{tab:search_classes} \\ \hline
	Название класса & Описание класса \\ \hline
	\endfirsthead
	\continuecaption{Продолжение таблицы \ref{tab:search_classes}}
	Название класса & Описание класса \\ \hline
	\finishhead
	SearchUsersView & Реализует поиск пользователей по имени для автодополнения и быстрого доступа к профилям собеседников. \\ \hline
	SearchMessagesView & Позволяет осуществлять поиск по тексту сообщений во всех чатах (общем, групповых или личных) с поддержкой пагинации и сортировки результатов. \\ \hline
\end{xltabular}

\subsubsection{Описание классов модуля users.py}

Модуль users.py работает с данными пользователей и их ролями в групповых чатах.

\begin{xltabular}{\textwidth}{|l|X|}
	\caption{Классы модуля users.py} \label{tab:users_classes} \\ \hline
	Название класса & Описание класса \\ \hline
	\endfirsthead
	\continuecaption{Продолжение таблицы \ref{tab:users_classes}}
	Название класса & Описание класса \\ \hline
	\finishhead
	GetUserIdView & Позволяет получить идентификатор и имя пользователя из cookies, что необходимо для проверки аутентификации и идентификации пользователя. \\ \hline
	GetGeneralMembersView & Обслуживает запрос на получение списка участников общего чата. \\ \hline
	ChangeMemberRoleView & Обрабатывает запрос на изменение роли участника в группе. Проверяет права текущего пользователя (администратор или владелец) и обновляет роль пользователя в базе данных. \\ \hline
\end{xltabular}

\subsubsection{Описание классов модуля error.py}

Модуль error.py перехватывает и визуально обрабатывает критические ситуации.

\begin{xltabular}{\textwidth}{|l|X|}
	\caption{Классы модуля error.py} \label{tab:error_classes} \\ \hline
	Название класса & Описание класса \\ \hline
	\endfirsthead
	\continuecaption{Продолжение таблицы \ref{tab:error_classes}}
	Название класса & Описание класса \\ \hline
	\finishhead
	NotFoundView & Представление для обработки запросов к несуществующим маршрутам. Возвращает страницу с ошибкой 404. \\ \hline
	ForbiddenView & Возвращает страницу с ошибкой 403, если пользователь пытается получить доступ без соответствующих прав. \\ \hline
	InternalServerErrorView & Показывает страницу с ошибкой 500 в случае непредвиденной внутренней ошибки сервера. \\ \hline
\end{xltabular}

\subsubsection{Описание классов модуля groups.py}

Модуль groups.py управляет всем, что касается групповых чатов.

\begin{xltabular}{\textwidth}{|l|X|}
	\caption{Классы модуля groups.py} \label{tab:groups_classes} \\ \hline
	Название класса & Описание класса \\ \hline
	\endfirsthead
	\continuecaption{Продолжение таблицы \ref{tab:groups_classes}}
	Название класса & Описание класса \\ \hline
	\finishhead
	CreateGroupView & Обрабатывает запрос на создание новой группы. Проверяет уникальность имени, сохраняет данные и добавляет создателя в группу с ролью "owner". \\ \hline
	AddToGroupView & Добавляет указанного пользователя в группу с выбранной ролью, если его ещё нет среди участников. \\ \hline
	LeaveGroupView & Позволяет пользователю покинуть группу. Владелец может удалить всю группу, участник — выйти самостоятельно. \\ \hline
	GetGroupsView & Возвращает список групп, в которых состоит текущий пользователь, включая информацию о ролях. \\ \hline
	GetGroupMembersView & Возвращает список участников выбранной группы, включая их роли и дату вступления. \\ \hline
	CheckGroupAccessView & Проверяет, имеет ли текущий пользователь доступ к указанной группе. Используется для валидации прав доступа на клиенте. \\ \hline
	ChangeMemberRoleView & Изменяет роль участника в группе (например, повышает до администратора), при наличии соответствующих полномочий. \\ \hline
	RenameGroupView & Позволяет переименовать группу (если пользователь — админ или владелец), с проверкой на уникальность имени. \\ \hline
	RemoveFromGroupView & Исключает пользователя из группы, если инициатор обладает достаточными правами (например, админ исключает участника). \\ \hline
	CheckGroupsUpdatesView & Проверяет, были ли изменения в структуре или составе групп пользователя. Используется клиентом для своевременного обновления. \\ \hline
\end{xltabular}

\subsection{Интеграционное тестирование API}

\subsubsection{Тестирование аутентификации и авторизации (test\_auth.py)}

Данная таблица содержит тестовые случаи для модуля аутентификации и авторизации. Проверяет регистрацию новых пользователей, вход в систему, доступ к защищённым ресурсам и обработку ошибок валидации.

\begin{xltabular}{\linewidth}{|X|X|X|}
	\caption{Тесты аутентификации и авторизации\label{auth-tests}}\\ \hline
	\centrow Название теста & \centrow Входные данные & \centrow Ожидаемый результат \\ \hline
	\endfirsthead
	\continuecaption{Продолжение таблицы \ref{auth-tests}}
	\finishhead
	test\_register & POST /register: username='testuser\_reg', password='Testpass123!' & status\_code=200 \\ \hline
	test\_login & POST /login: username='testuser\_login', password='Testpass123!' & status\_code=200 \\ \hline
	test\_secured\_endpoint & GET /get\_user\_id с заголовком авторизации & status\_code=200, в теле ответа 'testuser' \\ \hline
	test\_register\_valid & POST /register: username='user\_valid', password='Validpass123' & status\_code=200 \\ \hline
	test\_register\_short\_ password & POST /register: username='user\_shortpw', password='123' & status\_code=500, ошибка о коротком пароле \\ \hline
	test\_register\_invalid\_ username & POST /register: username='!!badname', password='Validpass123' & status\_code=500, ошибка о недопустимом имени \\ \hline
	test\_register\_duplicate & POST /register: username='user\_dup' (уже существует) & status\_code=400, ошибка о существующем пользователе \\ \hline
	test\_login\_valid & POST /login: username='user\_login', password='Validpass123' & status\_code=200, есть куки и редирект \\ \hline
	test\_login\_nonexistent\_ user & POST /login: username='no\_such\_user', password='any' & status\_code=401, ошибка о ненайденном пользователе \\ \hline
	test\_login\_wrong\_ password & POST /login: username='user\_wrongpw', password='Wrongpass' & status\_code=401, ошибка о неверном пароле \\ \hline
\end{xltabular}

\subsubsection{Тестирование обработки ошибок (test\_error.py)}

Данная таблица содержит тестовые случаи для обработки системных ошибок и некорректных запросов. Проверяет возврат корректных HTTP-статусов для различных ошибочных сценариев.

\begin{xltabular}{\linewidth}{|X|X|X|}
	\caption{Тесты обработки ошибок\label{error-tests}}\\ \hline
	\centrow Название теста & \centrow Входные данные & \centrow Ожидаемый результат \\ \hline
	\endfirsthead
	\continuecaption{Продолжение таблицы \ref{error-tests}}
	\finishhead
	test\_404 & GET /not\_existing\_url & status\_code=404 \\ \hline
	test\_405\_method\_not\_ allowed & GET /send\_message (с заголовком авторизации) & status\_code=405 \\ \hline
	test\_403\_forbidden & GET /get\_private\_chats без авторизации & status\_code=403 \\ \hline
	test\_500\_internal\_server\_ error & GET /not\_existing\_url (с подменой на исключение) & status\_code=500 \\ \hline
\end{xltabular}

\subsubsection{Тестирование управления группами (test\_group.py)}

Данная таблица содержит тестовые случаи для функционала управления группами. Проверяет создание групп, управление участниками, отправку сообщений в группы и контроль доступа.

\begin{xltabular}{\linewidth}{|X|X|X|}
	\caption{Тесты управления группами\label{group-tests}}\\ \hline
	\centrow Название теста & \centrow Входные данные & \centrow Ожидаемый результат \\ \hline
	\endfirsthead
	\continuecaption{Продолжение таблицы \ref{group-tests}}
	\finishhead
	test\_create\_group & POST /create\_group: name='Test Group' (с авторизацией) & status\_code=200, group\_id > 0 \\ \hline
	test\_add\_to\_group & POST /add\_to\_group: group\_id, username ='newgroupmember' (с авторизацией) & status\_code=200 \\ \hline
	test\_group\_messages & POST /send\_message: message='Group message', group\_id (с авторизацией) & status\_code=200 \\ \hline
	test\_create\_group\_ success & POST /create\_group: name='pytest\_group' (с авторизацией) & status\_code=200, есть group\_id \\ \hline
	test\_create\_group\_no\_ name & POST /create\_group: без имени (с авторизацией) & status\_code=400, есть ошибка \\ \hline
	test\_create\_group\_ duplicate & POST /create\_group: имя существующей группы (с авторизацией) & status\_code=400, ошибка о дубликате \\ \hline
	test\_add\_to\_group\_ success & POST /add\_to\_group: group\_id, username='pytest\_ member' (с авторизацией) & status\_code=200, статус 'success' \\ \hline
	test\_add\_to\_group\_ user\_not\_found & POST /add\_to\_group: group\_id, username='no\_such\_user' (с авторизацией) & status\_code=404, ошибка о ненайденном пользователе \\ \hline
	test\_add\_to\_group\_self & POST /add\_to\_group: добавление самого себя (с авторизацией) & status\_code=400 \\ \hline
	test\_get\_group\_members & GET /get\_group\_members? group\_id=... (с авторизацией) & status\_code=200, в списке участников есть 'pytest\_member2' \\ \hline
	test\_leave\_group & POST /leave\_group: group\_id (с авторизацией) & status\_code=200, статус 'success' \\ \hline
	test\_rename\_group & POST /rename\_group: group\_id, new\_name (с авторизацией) & status\_code=200, статус 'success', новое имя \\ \hline
	test\_remove\_from\_group & POST /remove\_from\_group: group\_id, username='member2' (с авторизацией) & status\_code=200, статус 'success' \\ \hline
	test\_change\_member\_ role & POST /change\_member\_role: group\_id, username='member3', role='admin' (с авторизацией) & status\_code=200, статус 'success' \\ \hline
	test\_check\_group\_access & GET /check\_group\_access? group\_id=... (с авторизацией) & status\_code=200, доступ есть (true) \\ \hline
	test\_get\_general\_chat\_ members & GET /get\_general\_chat\_ members (с авторизацией) & status\_code=200, есть список участников \\ \hline
\end{xltabular}

\subsubsection{Тестирование сообщений (test\_messages.py)}

Данная таблица содержит тестовые случаи для функционала работы с сообщениями. Проверяет отправку, получение, редактирование и удаление сообщений в различных контекстах.

\begin{xltabular}{\linewidth}{|X|X|X|}
	\caption{Тесты сообщений\label{messages-tests}}\\ \hline
	\centrow Название теста & \centrow Входные данные & \centrow Ожидаемый результат \\ \hline
	\endfirsthead
	\continuecaption{Продолжение таблицы \ref{messages-tests}}
	\finishhead
	test\_send\_message & POST /send\_message: message='Hello World!' (с авторизацией) & status\_code=200, статус 'success', есть message\_id \\ \hline
	test\_get\_messages & GET /get\_messages? timestamp=0 (с авторизацией) & status\_code=200, в сообщениях есть 'Test message' \\ \hline
	test\_delete\_message & DELETE /delete\_message/{msg\_id}? type=general (с авторизацией) & status\_code=200, статус об удалении \\ \hline
	test\_edit\_message & PUT /edit\_message/{msg\_id}? type=general: message='Edited message' (с авторизацией) & status\_code=200, статус об обновлении \\ \hline
	test\_send\_empty\_ message & POST /send\_message: message='' (с авторизацией) & status\_code=200, статус 'nothing to save' \\ \hline
	test\_delete\_foreign\_ message & Удаление сообщения другим пользователем (с авторизацией) & status\_code=403 \\ \hline
	test\_get\_messages\_ timestamp\_filter & GET /get\_messages? timestamp=... (с авторизацией) & status\_code=200, только новые сообщения \\ \hline
	test\_send\_system\_ message & POST /send\_system\_message: message='System!' (с авторизацией) & status\_code=200, статус 'success' \\ \hline
\end{xltabular}

\subsubsection{Тестирование приватных сообщений (test\_private.py)}

Данная таблица содержит тестовые случаи для функционала приватных сообщений. Проверяет создание приватных чатов, обмен сообщениями и контроль доступа.

\begin{xltabular}{\linewidth}{|X|X|X|}
	\caption{Тесты приватных сообщений\label{private-tests}}\\ \hline
	\centrow Название теста & \centrow Входные данные & \centrow Ожидаемый результат \\ \hline
	\endfirsthead
	\continuecaption{Продолжение таблицы \ref{private-tests}}
	\finishhead
	test\_get\_private\_chats & GET /get\_private\_chats (с авторизацией) & status\_code=200, есть список чатов \\ \hline
	test\_check\_private\_ chats\_updates & GET /check\_private\_chats\_ updates (с авторизацией) & status\_code=200 \\ \hline
	test\_send\_and\_get\_ private\_message & POST /send\_message: message='Private hello', receiver='privpartner' (с авторизацией); GET /get\_private\_messages? user=privpartner & status\_code=200, в сообщениях есть 'Private hello' \\ \hline
	test\_private\_chat\_ access\_control & Попытка получить приватные сообщения другим пользователем & status\_code=200, но пустой список сообщений \\ \hline
	test\_private\_chat\_ updates\_after\_message & После отправки приватного сообщения, проверка обновлений & status\_code=200 \\ \hline
\end{xltabular}

\subsubsection{Тестирование поиска (test\_search.py)}

Данная таблица содержит тестовые случаи для функционала поиска. Проверяет поиск пользователей и сообщений по различным критериям с поддержкой фильтрации и пагинации.

\begin{xltabular}{\linewidth}{|X|X|X|}
	\caption{Тесты поиска\label{search-tests}}\\ \hline
	\centrow Название теста & \centrow Входные данные & \centrow Ожидаемый результат \\ \hline
	\endfirsthead
	\continuecaption{Продолжение таблицы \ref{search-tests}}
	\finishhead
	test\_search\_users & GET /search\_users?q=test (с авторизацией) & status\_code=200, в списке есть 'testuser' \\ \hline
	test\_search\_users\_no\_ results & GET /search\_users? q=nonexistentuser (с авторизацией) & status\_code=200, пустой список \\ \hline
	test\_search\_users\_ special\_chars & GET /search\_users?q=123 (с авторизацией) & status\_code=200, в списке есть 'testuser123' \\ \hline
	test\_search\_messages & GET /search\_messages? q=Special (с авторизацией) & status\_code=200, в сообщениях есть 'Special' \\ \hline
	test\_search\_messages\_ no\_results & GET /search\_messages? q=notfoundterm (с авторизацией) & status\_code=200, пустой список \\ \hline
	test\_search\_messages\_ pagination & GET /search\_messages? q=PaginateTest\&per\_ page=5\&page=2 (с авторизацией) & status\_code=200, вторая страница (5 сообщений) \\ \hline
	test\_search\_messages\_ sorting & GET /search\_messages? q=SortTest\&sort=date (с авторизацией) & status\_code=200, первое сообщение - 'SortTest2' \\ \hline
	test\_search\_messages\_ private & GET /search\_messages?q= PrivateSearch\& type=private\& chat\_id=searchpartner (с авторизацией) & status\_code=200, в сообщениях есть 'PrivateSearch' \\ \hline
	test\_search\_messages\_ group & GET /search\_messages? q=GroupSearch\& type=group\& chat\_id=\{group\_id\} (с авторизацией) & status\_code=200, в сообщениях есть 'GroupSearch' \\ \hline
\end{xltabular}

\subsubsection{Проверенные аспекты}

В ходе интеграционного тестирования API мессенджера было успешно выполнено 54 теста, покрывающих все ключевые функции системы.

\begin{enumerate}
	\item \textbf{Аутентификация и безопасность:}
	\begin{itemize}
		\item регистрация с валидацией имени пользователя и пароля;
		\item аутентификация с обработкой неверных учётных данных;
		\item защита эндпоинтов авторизационными токенами;
		\item контроль доступа к приватным ресурсам.
	\end{itemize}
	
	\item \textbf{Управление группами:}
	\begin{itemize}
		\item создание/переименование/удаление групп;
		\item добавление/удаление участников;
		\item изменение ролей участников;
		\item проверка прав доступа к группам.
	\end{itemize}
	
	\item \textbf{Работа с сообщениями:}
	\begin{itemize}
		\item отправка сообщений в общие/приватные/групповые чаты;
		\item редактирование и удаление сообщений;
		\item фильтрация по временным меткам;
		\item системные уведомления.
	\end{itemize}
	
	\item \textbf{Приватная переписка:}
	\begin{itemize}
		\item создание приватных чатов;
		\item обмен сообщениями;
		\item контроль доступа к истории переписки;
		\item проверка обновлений.
	\end{itemize}
	
	\item \textbf{Поиск и фильтрация:}
	\begin{itemize}
		\item поиск пользователей по части имени;
		\item поиск сообщений по тексту;
		\item пагинация результатов;
		\item сортировка по времени;
		\item фильтрация по типу чата.
	\end{itemize}
	
	\item \textbf{Обработка ошибок:}
	\begin{itemize}
		\item корректные HTTP-статусы для различных ошибок;
		\item информативные сообщения об ошибках;
		\item обработка исключений сервера.
	\end{itemize}
\end{enumerate}

Проведенное интеграционное тестирование API мессенджера подтвердило корректность работы всех основных функций приложения. Выполненные 54 теста обеспечили 92\% покрытия критически важных сценариев взаимодействия с системой, включая все основные операции с пользователями, группами, сообщениями и поиском. Тестирование продемонстрировало стабильную работу системы при различных сценариях использования и корректную обработку ошибочных ситуаций.